\section{Conclusion: Boundary Conditions of RSI and Affective-Memory Mechanisms}
\label{sec:conclusion}
% 【中文】\section{结论：RSI 与情绪记忆机制的边界条件}

We pre-registered 3 core hypotheses (H1--H3) and 7 sub-hypotheses on whether governed long-term memory + RSI and three affective-memory mechanisms yield transferable gains on long-chain EDA flow. M0 (\MzeroPass{}) anchors the difficulty; other arms are evaluated after \PreregDate.
% 【中文】
% 我们预注册了 3 条核心假设（H1--H3）与 7 条子假设，探讨受治理的长期记忆 + RSI 与三个情绪记忆机制
% 是否在长链 EDA 流程上带来可迁移增益。M0（\MzeroPass{}）锚定难度；其他臂在 \PreregDate 之后评测。

\textbf{Honest route.} If H1/H2 are falsified, this paper reports boundary conditions of RSI / affective heuristics in long-chain flow---both publishable as they delimit where bio-inspired mechanisms apply. If supported, it provides the first pre-registered evidence that affective-memory mechanisms bring mechanistic, not merely functional, gains to long-chain EDA orchestration.
% 【中文】
% \textbf{诚实路线。} 若 H1/H2 被证伪，本文报告长链流程中 RSI / 情绪启发式的边界条件——两者均可发表，
% 因为它们划定了仿生机制的适用范围。若获支持，本文提供首个预注册证据：情绪记忆机制为长链 EDA 编排
% 带来机制性而非纯功能性的增益。

\textbf{Future work.} Multi-modal grounding; longer chains ($>$50 calls); finer valence dimensions; and whether H3 reverses at extreme chain lengths.
% 【中文】
% \textbf{未来工作。} 多模态接地；更长链（$>$50 次调用）；更细的效价维度；以及 H3 是否在极长链长度下反转。
